% Part: first-order-logic
% Chapter: completeness
% Section: lindenbaums-lemma

\documentclass[../../../include/open-logic-section]{subfiles}

\begin{document}

\iftag{FOL}
      {\olfileid[fr]{fol}{com}{lin}}
      {\olfileid[fr]{pl}{com}{lin}}

\olsection{Lemme de Lindenbaum}

\begin{explain}
Nous allons montrer que tout ensemble cohérent \iftag{FOL}{d'!!{sentence}s}{de !!{formula}s} est
contenu dans un ensemble à la fois cohérent et !!{complete}. La
démonstration ajoute successivement des \iftag{FOL}{énoncés}{formules}, en maintenant la
cohérence à chaque étape. Elle procède de sorte que, pour \iftag{FOL}{tout !!{sentence}}{toute !!{formula}}~$!A$, soit $!A$, soit $\lnot !A$ soit ajouté à une
étape. La réunion de toutes les étapes contient donc $!A$ ou sa
négation et est ainsi complète. Elle reste cohérente, car aucune
étape n'introduit d'incohérence.
\end{explain}

\begin{lem}[Lemme de Lindenbaum]
\ollabel{lem:lindenbaum} Tout ensemble cohérent~$\Gamma$ dans un
langage~$\Lang{L}$ peut s'étendre en un ensemble cohérent et
!!{complete}~$\Gamma^*$.
\end{lem}

\begin{proof}
Soit $\Gamma$ cohérent. Dans le langage au plus dénombrable fixé ici,
énumérons \iftag{FOL}{tous les !!{sentence}s}{toutes les !!{formula}s} de~$\Lang L$ sous la forme
$!A_0,!A_1,\dots{}$. Posons $\Gamma_0=\Gamma$, puis
\[
\Gamma_{n+1} =
\begin{cases}
\Gamma_n \cup \{ !A_n \} & \textrm{si $\Gamma_n \cup \{!A_n\}$ est
  cohérent ;} \\
\Gamma_n \cup \{ \lnot !A_n \} & \textrm{sinon.}
\end{cases}
\]
Enfin, posons $\Gamma^*=\bigcup_{n\geq 0}\Gamma_n$.

Chaque $\Gamma_n$ est cohérent. En effet, $\Gamma_0$ l'est par
hypothèse. Si $\Gamma_{n+1}=\Gamma_n\cup\{!A_n\}$, cette réunion
est cohérente par le choix de la première branche. Sinon,
$\Gamma_{n+1}=\Gamma_n\cup\{\lnot !A_n\}$. Vérifions que cette
dernière réunion est cohérente. Si elle ne l'était pas,
\emph{les deux} ensembles $\Gamma_n\cup\{!A_n\}$ et
$\Gamma_n\cup\{\lnot !A_n\}$ seraient incohérents.
D'après les
\iftag{FOL}{%
  \tagrefs{prfAX/{fol:axd:prv:prop:provability-exhaustive},
    prfSC/{fol:seq:prv:prop:provability-exhaustive},
    prfND/{fol:ntd:prv:prop:provability-exhaustive},
    prfTab/{fol:tab:prv:prop:provability-exhaustive}}}{%
  \tagrefs{prfAX/{pl:axd:prv:prop:provability-exhaustive},
    prfSC/{pl:seq:prv:prop:provability-exhaustive},
    prfND/{pl:ntd:prv:prop:provability-exhaustive},
    prfTab/{pl:tab:prv:prop:provability-exhaustive}}},
$\Gamma_n$ serait alors incohérent, contrairement à l'hypothèse de
récurrence.

Pour tout~$n$ et tout $i<n$, on a $\Gamma_i\subseteq\Gamma_n$.
Montrons-le par récurrence sur~$n$. Pour $n=0$, il n'y a aucun
$i<0$ : l'énoncé est donc immédiat. Supposons-le vrai pour~$n$.
Par construction, $\Gamma_{n+1}$ vaut
$\Gamma_n\cup\{!A_n\}$ ou
$\Gamma_n\cup\{\lnot !A_n\}$ ; ainsi,
$\Gamma_n\subseteq\Gamma_{n+1}$. Si $i<n+1$, alors
$\Gamma_i\subseteq\Gamma_n$ par hypothèse de récurrence quand
$i<n$, et par réflexivité de~$\subseteq$ quand $i=n$.
La transitivité de~$\subseteq$ donne donc
$\Gamma_i\subseteq\Gamma_{n+1}$.

Il en résulte que $\Gamma^*$ est cohérent. Soit, en effet,
$\Gamma'\subseteq\Gamma^*$ fini. Si $\Gamma'$ est vide, prenons
$n=0$. Sinon, pour chaque $!B\in\Gamma'$ choisissons un indice~$i$
tel que $!B\in\Gamma_i$, et prenons pour~$n$ le plus grand de ces
indices. Puisque $\Gamma_i\subseteq\Gamma_n$ dès que $i\le n$,
tous les éléments de~$\Gamma'$ appartiennent à~$\Gamma_n$.
Ainsi, $\Gamma'\subseteq\Gamma_n$ et $\Gamma'$ est cohérent.
Tout sous-ensemble fini de~$\Gamma^*$ est donc cohérent. Par les
\iftag{FOL}{%
  \tagrefs{prfAX/{fol:axd:ptn:prop:proves-compact},
    prfSC/{fol:seq:ptn:prop:proves-compact},
    prfND/{fol:ntd:ptn:prop:proves-compact},
    prfTab/{fol:tab:ptn:prop:proves-compact}}}{%
  \tagrefs{prfAX/{pl:axd:ptn:prop:proves-compact},
    prfSC/{pl:seq:ptn:prop:proves-compact},
    prfND/{pl:ntd:ptn:prop:proves-compact},
    prfTab/{pl:tab:ptn:prop:proves-compact}}},
$\Gamma^*$ est cohérent.

Enfin, \iftag{FOL}{tout !!{sentence}}{toute !!{formula}} du langage apparaît dans l'énumération
utilisée. Si $!A_n\notin\Gamma^*$, c'est que
$\Gamma_n\cup\{!A_n\}$ était incohérent. La seconde branche a
donc ajouté $\lnot !A_n$ à~$\Gamma^*$. Ainsi, pour \iftag{FOL}{tout !!{sentence}}{toute !!{formula}}~$!A_n$, $!A_n\in\Gamma^*$ ou
$\lnot !A_n\in\Gamma^*$ : l'ensemble $\Gamma^*$ est
!!{complete}.
\end{proof}

\end{document}
